\chapter{Metodología}
\label{cap:metodologia}

\section{Tipo de investigación}
\label{sec:tipo-investigacion}

El proyecto combina investigación aplicada y desarrollo tecnológico: a partir del planteamiento del anteproyecto se define un diseño formal del protocolo (Incremento~1), se implementa en software (Incremento~2 y siguientes) y se valida mediante experimentos controlados en laboratorio con dispositivos Android comerciales.

\section{Enfoque incremental}
\label{sec:enfoque-incremental}

La metodología sigue entregas incrementales trazables:

\begin{enumerate}[label=\textbf{I\arabic*.}]
  \item \textbf{Incremento 1 --- Diseño formal:} arquitectura, requisitos, wire format TCP, máquinas de estado, matriz de validación (\texttt{docs/protocol/}, versión spec 0.1.0-pre).
  \item \textbf{Incremento 2 --- Capas formales:} modelo OSI emulado, L2 (enlace TCP dual), L3 (red), envoltorio de paquetes (spec 0.2.1-draft).
  \item \textbf{Implementación por fases A--D:} enlace $\rightarrow$ tabla de rutas $\rightarrow$ reenvío multisalto $\rightarrow$ endurecimiento.
  \item \textbf{Validación experimental:} experimentos EXP-01 a EXP-06 con criterios de aceptación medibles.
\end{enumerate}

\section{Fases de implementación}
\label{sec:fases-implementacion}

La Tabla~\ref{tab:fases} resume las fases de la hoja de ruta (\texttt{17-implementation-roadmap.md}).

\begin{table}[htbp]
  \centering
  \caption{Fases de implementación lib-pct-core}
  \label{tab:fases}
  \begin{tabular}{@{}ll@{}}
    \toprule
    \textbf{Fase} & \textbf{Objetivo} \\
    \midrule
    A & Capa de enlace formal: TCP :8765/:8766, NeighborRegistry \\
    B & Tabla de rutas local y TOPO\_UPDATE \\
    C & Reenvío multisalto (ForwardWorker) en canal datos \\
    D & NODE\_DOWN, rutas STALE, persistencia UUID, \texttt{\_pct-seek} \\
    \bottomrule
  \end{tabular}
\end{table}

\section{Matriz de validación experimental}
\label{sec:matriz-validacion}

\subsection{Hipótesis}

\begin{table}[htbp]
  \centering
  \caption{Hipótesis principales}
  \label{tab:hipotesis}
  \begin{tabular}{@{}cl@{}}
    \toprule
    \textbf{ID} & \textbf{Hipótesis} \\
    \midrule
    H1 & GO propio + STA legacy upstream simultáneos en clase $\alpha$ \\
    H2 & Dos GO interconectados sin \texttt{connect()} \\
    H3 & Cadena de tres nodos con reenvío TCP \\
    H4 & Unión tardía sin reinicio de SSID \\
    H5 & Reconexión mantiene \texttt{session\_epoch} \\
    \bottomrule
  \end{tabular}
\end{table}

\subsection{Experimentos}

\textbf{EXP-01 (gate crítico):} interconexión de dos GO. Criterio: hijo mantiene GO propio tras STA; HELLO + JOIN\_COMMIT exitosos; $T_{\mathrm{join}} < 15$~s.

\textbf{EXP-02:} DNS-SD con GO local sin Group Client.

\textbf{EXP-03:} cadena ROOT $\leftarrow$ BRIDGE $\leftarrow$ LEAF; DATA de C a A vía reenvío en B.

\textbf{EXP-04:} unión tardía de nodo ISLAND a red operativa.

\textbf{EXP-05:} caída de BRIDGE y reconvergencia de sub-árbol.

\textbf{EXP-06:} LEAF anuncia \texttt{\_pct-seek} y recibe JOIN\_OFFER.

\section{Entorno de laboratorio}
\label{sec:entorno-lab}

\begin{itemize}
  \item \textbf{Dispositivos:} mínimo cuatro smartphones clase~$\alpha$ (p.\ ej.\ Samsung Galaxy A24, Vivo Y21s).
  \item \textbf{Herramientas:} Android Studio, Gradle, ADB, logcat (\texttt{PctMesh}).
  \item \textbf{Software:} app demo \texttt{pctmesh}, biblioteca \texttt{co.uan.pct:core}.
  \item \textbf{Red:} entorno indoor sin AP externo; grupos P2P aislados.
\end{itemize}

\section{Técnicas de recolección de datos}
\label{sec:recoleccion-datos}

\begin{enumerate}
  \item Logs estructurados en tag \texttt{PctMesh} (bootstrap, L2, DNS-SD).
  \item Snapshots de UI debug: fase, vecinos, rutas, diagnóstico DNS-SD.
  \item Medición de $T_{\mathrm{join}}$, $T_{\mathrm{reconfig}}$ con marca temporal en logs.
  \item Tests unitarios JUnit: codec TCP, parseo UUID, utilidades de red.
\end{enumerate}

\section{Síntesis del capítulo}
\label{sec:sintesis-metodologia}

La metodología garantiza trazabilidad desde requisitos hasta código y experimentos. EXP-01 actúa como compuerta: su fallo obliga a reclasificar hardware antes de pruebas multisalto.
